\section{Introduction}
\label{sec:introduction}


Ovarian cancer is one of the cancers with the highest incidence and mortality rates among women worldwide~\cite{sung2021global}. 
The general treatment is cytoreductive surgery; however, residual disease from an incompletely successful surgery may cause extensive metastasis, increasing the risk of postoperative recurrence and mortality~\cite{chen2021ct}.
Therefore, it is crucial to accurately predict the residual disease after cytoreductive surgery to make clinical surgical decisions.
In this paper, we refer to the objective as \emph{residual disease prediction of ovarian cancer.}

%For the accurate prediction of residual disease of ovarian cancer, large efforts have been made by researchers from various areas, such as medical and computer science.
As far as we are concerned, large efforts have been made by researchers from various fields, such as medical and computer science, to accurately predict the residual disease of ovarian cancer.
In earlier stages, laparoscopy, a kind of invasive examination, was applied to provide a visual and systematic view of the whole abdomen \cite{testa2012ultrasound, van2019laparoscopy}.
Later, Li et al. \cite{li2021noninvasive} employed feature engineering and statistical method to analyze clinical and medical imaging data obtained by noninvasive examination.
With the successful accomplishment of deep learning, significant progress has been made in medical image analysis.
Wang et al. \cite{wang2019deep}, proposed to leverage a CNN-based framework to extract deep image features for recurrence prediction of ovarian cancer.
\begin{figure}
    \centering
    \includegraphics[height=45mm, width=65mm]{figure/MuVAL_sample.pdf}
    % \vspace{-4mm}
    \caption{Multi-view of a 3D CT image.}
    \label{f:MuVAL_sample}
    % \vspace{-4mm}
\end{figure}


However, traditional invasive methods sometimes naturally carry an additional risk of metastasis.
Furthermore, feature engineering and statistical methods often require manual selection and design of features by knowledgeable medical professionals, which is both time-consuming and labor-intensive.
%while feature engineering and statistic methods often require manual selection and design of features by experienced doctors, which is time- and labor-consuming.
In addition, despite the remarkable progress of deep-learning technology in medical image analysis,
most of the current deep-learning methods underestimate the importance of 3D image information contained in Computed Tomography (CT) or other medical images, which may bring a limited performance for residual disease prediction.

As a matter of fact, a 3D CT image of a patient's entire disease region is composed of a stack of $d$ 2D CT slices, each with a resolution of ($h, w$). 
According to the size of the disease region, $d$ is usually between 10 and 100.
%Containing the whole disease region of a patient, a 3D CT image consists of a stack of $d$ 2D CT slices, each with a resolution of ($h, w$). Depending on the size of the disease region, $d$ is typically between 10 and 100. 
% height of $h$ and a width of $w$.
%As shown in Fig~\ref{f:MuVAL_sample}, traditionally, the 2D CT slices are the transverse view (in green). 
Traditionally, the 2D CT slices used in the previous works only included the transverse view (in green), as shown in Fig~\ref{f:MuVAL_sample}.
However, multi-view features facilitate a more comprehensive and accurate observation of 3D images~\cite{wildeman20092d, wu2020deep}.
Thus, by fully utilizing multi-view features such as the coronal view (in blue) and the sagittal view (in red), more specific information about the disease region can be obtained. 
%%%%%%扩充挑战
Although the value of multi-view features has been proved, there are still many challenges in better utilizing multi-view features in 3D CT images. 
The difficulty of image representation is compounded by factors such as the shape, size, texture, invasion, and metastasis of the disease, all of which can affect the outcome of surgery. Additionally, the informativeness of 2D CT slices varies, with some slices reflecting the residual disease situation better than others. 
Further, differences in cancer subtypes and surgical teams introduce differences in surgical procedures, and specific surgical teams and cancer subtypes will limit the scale of available data.


To this end, in this paper, we propose a novel Multi-View Attention Learning~(MuVAL), an end-to-end deep-learning method for residual disease prediction of ovarian cancer. 
Specifically, we first obtain transverse, coronal, and sagittal views of 3D CT images to create multi-view images.
To better represent the image features, we leverage Squeeze-and-Excitation~(SE) \cite{he2016deep} module to help find the more relevant slices in each view.  
In the backbone, we apply Med3d \cite{chen2019med3d}, a pre-trained ResNet \cite{he2016deep} for 3D medical image, to deal with the small-scale dataset.
Extensive experiments and multiple metrics demonstrate the superiority of our proposed method.
As an emphasis, the main contributions of our work can be concluded as follow: 
\begin{itemize}
    \item We observe the great potential of utilizing multi-view of 3D CT images from transverse, coronal, and sagittal views, and propose to leverage the multi-view features to enhance the learning for residual disease prediction of ovarian cancer.
    \item We propose a novel Multi-View Attention Learning (MuVAL) method for residual disease prediction of ovarian cancer, which focuses on the comprehensive learning of 3D CT images in a multi-view manner.
    \item Extensive experiments and multiple metrics demonstrate the superiority and rationality of our proposed method compared with the baseline methods.
\end{itemize}